% !TeX root = Report.tex
\chapter{Описание алгоритмов и схема оценивания}

В данной главе описывается математическая постановка задачи и вводятся
алгоритмы, используемые далее в численном исследовании.

\section{Модель изображения и параметризация прямой}

Рассматривается изображение в градациях серого, представленное матрицей
\[
\mathbf{M} = (M_{ij}) \in \mathbb{R}^{N_{\mathrm{row}} \times N_{\mathrm{col}}}.
\]
Здесь $N_{\mathrm{row}}$ --- число строк изображения, а $N_{\mathrm{col}}$ ---
число столбцов. Предполагается, что наблюдаемое изображение раскладывается в
сумму нескольких компонент сигнала и шумовой компоненты:
\[
\mathbf{M} = \sum_{\ell=1}^{q} \mathbf{S}^{(\ell)} + \mathbf{R},
\]
где $q$ --- число прямых на изображении, $\mathbf{S}^{(\ell)}$ --- матрица
сигнала, соответствующая $\ell$-й прямой, а элементы шумовой матрицы
$\mathbf{R}$ независимы и имеют распределение $\mathcal{N}(0,\sigma^2)$.

Пусть $\ell$-я прямая задается параметрами $a_\ell$ и $b_\ell$. Тогда ей
соответствует множество
\[
\Gamma_\ell
=
\left\{
(i,j)\in\{1,\dots,N_{\mathrm{row}}\}\times\{1,\dots,N_{\mathrm{col}}\}
\colon
i = [a_\ell j+b_\ell]
\right\},
\]
где $[\cdot]$ обозначает целую часть. Матрица сигнала для этой прямой имеет вид
\[
S^{(\ell)}_{ij}
=
c_\ell \mathbf{1}\{(i,j)\in\Gamma_\ell\}.
\]
Значения вне $\Gamma_\ell$ равны нулю, а внутри равны числу $c_\ell>0$,
которое мы будем называть интенсивностью. Если несколько прямых проходят через
один и тот же пиксель, итоговая яркость в этой точке определяется суммой
соответствующих интенсивностей.

\section{Преобразование Хафа}

Для описания прямых используется параметризация
\[
\rho = x \cos \theta + y \sin \theta,
\]
где $\rho$ --- расстояние от начала координат до прямой, а $\theta$ --- угол
нормали к прямой \cite{duda1972hough}. Такая запись позволяет рассматривать
как наклонные, так и вертикальные прямые.

\begin{figure}[H]
	\centering
	\includegraphics[width=0.5\textwidth]{rt_params.pdf}
	\caption{Параметризация прямой через $(\rho,\theta)$}
	\label{fig:rt-params}
\end{figure}

Пусть после предобработки получена матрица $\widetilde{\mathbf{M}}$. Для
заданного порога $\varepsilon>0$ рассмотрим множество
\[
E_{\varepsilon}(\widetilde{\mathbf{M}})
=
\{(i,j)\colon \widetilde{M}_{ij}>\varepsilon\}.
\]
Будем называть $E_{\varepsilon}(\widetilde{\mathbf{M}})$ активным множеством,
а его элементы --- активными элементами матрицы.
На рис.~\ref{fig:one-line-example} показан пример бинарного изображения, в
котором активное множество соответствует прямой $y=x-10$.

\begin{figure}[H]
	\centering
	\includegraphics[width=0.44\textwidth]{chapter1_one_line_clean.png}
	\caption{Пример матрицы изображения с одной прямой $y=x-10$}
	\label{fig:one-line-example}
\end{figure}

Для фиксированной активной точки $(x_0,y_0)$ все прямые, проходящие через нее,
имеют параметры $(\rho,\theta)$, удовлетворяющие уравнению
\[
\rho(\theta; x_0, y_0) = \rho(\theta)=x_0\cos\theta+y_0\sin\theta.
\]
Если несколько активных точек лежат на одной прямой, то соответствующие кривые
в пространстве параметров имеют общую точку.

В рассматриваемой постановке задачи пространство параметров заменяется сеткой.
Параметр $\rho$ может принимать значения от $-\rho_{\max}$ до $\rho_{\max}$,
где
\[
\rho_{\max}
=
\left\lceil
\sqrt{N_{\mathrm{row}}^2+N_{\mathrm{col}}^2}
\right\rceil.
\]
При заданных шагах $h_\rho$ и $h_\theta$ дискретизации параметры принимают значения
\[
\rho_u = -\rho_{\max} + u h_\rho,
\qquad
u=0,\dots,\left\lfloor\frac{2\rho_{\max}}{h_\rho}\right\rfloor,
\]
\[
\theta_v = v h_\theta,
\qquad
v=0,\dots,\left\lfloor\frac{\pi}{h_\theta}\right\rfloor.
\]
То есть $\rho$ меняется от $-\rho_{\max}$ до последнего узла сетки, не
превосходящего $\rho_{\max}$, а $\theta$ --- от $0$ до последнего узла сетки,
не превосходящего $\pi$.

Аккумуляторный массив $A=(A_{uv})$ задается формулой
\[
A_{uv}
=
\sum_{(x_0,y_0)\in E_{\varepsilon}(\widetilde{\mathbf{M}})}
\mathbf{1}
\left\{
\left|
\rho_u - \rho(\theta_v;x_0,y_0)
\right|
\le \frac{h_\rho}{2}
\right\}.
\]
Другими словами, аккумуляторный массив кодирует пересечения кривых
$\rho(\theta)$. Большие значения $A_{uv}$ соответствуют тем параметрам
$(\rho_u,\theta_v)$, для которых большее число активных точек исходного
изображения лежит на одной прямой с такими параметрами. Поэтому локальные
максимумы аккумуляторного массива рассматриваются как оценки параметров прямых
на изображении:
\[
(\widehat{\rho},\widehat{\theta})
=
(\rho_{\widehat{u}},\theta_{\widehat{v}}),
\qquad
(\widehat{u},\widehat{v})
\in
\argmax_{u,v} A_{uv}.
\]
Алгоритм преобразования Хафа изложен ниже.

\begin{algorithm}[H]
	\caption{Алгоритм преобразования Хафа}
	\begin{algorithmic}[1]
		\small
		\Require матрица изображения $\widetilde{\mathbf{M}}\in\mathbb{R}^{N_{\mathrm{row}}\times N_{\mathrm{col}}}$, порог активности $\varepsilon>0$, шаги $h_\rho$, $h_\theta$
		\Ensure аккумуляторный массив $A$, оценка параметров прямой $(\widehat{\rho},\widehat{\theta})$
		\State $E_{\varepsilon}(\widetilde{\mathbf{M}})\gets\{(i,j):\widetilde{M}_{ij}>\varepsilon\}$
		\State $\rho_{\max}\gets\left\lceil\sqrt{N_{\mathrm{row}}^2+N_{\mathrm{col}}^2}\right\rceil$
		\State $U\gets\left\lfloor 2\rho_{\max}/h_\rho\right\rfloor$, \quad $V\gets\left\lfloor \pi/h_\theta\right\rfloor$
		\State $\rho_u\gets-\rho_{\max}+u h_\rho$, $u=0,\dots,U$; \quad $\theta_v\gets v h_\theta$, $v=0,\dots,V$; \quad $A_{uv}\gets 0$
		\For{каждого $(i,j)\in E_{\varepsilon}(\widetilde{\mathbf{M}})$ и каждого $v=0,\dots,V$}
			\State $\rho(\theta_v)\gets j\cos\theta_v+i\sin\theta_v$
			\State выбрать $u\in\{0,\dots,U\}: |\rho_u-\rho(\theta_v)|\le h_\rho/2$
			\State $A_{uv}\gets A_{uv}+1$
		\EndFor
		\State $(\widehat{u},\widehat{v})\in\argmax_{u,v}A_{uv}$
		\State положить $(\widehat{\rho},\widehat{\theta})\gets(\rho_{\widehat{u}},\theta_{\widehat{v}})$
		\State \Return $A$, $(\widehat{\rho},\widehat{\theta})$
	\end{algorithmic}
\end{algorithm}

На рис.~\ref{fig:accumulator} показан пример преобразования Хафа для
изображения с двумя прямыми $y=2x-1$ и $y=-x+101$. На аккумуляторном массиве
выделяются две области, в которых пересекается наибольшее число кривых
$\rho(\theta;x_k,y_k)$; положения этих областей определяют параметры
соответствующих прямых.

\begin{figure}[!htbp]
	\centering
	\begin{subfigure}{0.45\textwidth}
		\centering
		\includegraphics[width=\textwidth,height=0.30\textheight,keepaspectratio]{chapter1_two_lines_clean.png}
		\caption{Матрица изображения}
	\end{subfigure}
	\hfill
	\begin{subfigure}{0.45\textwidth}
		\centering
		\includegraphics[width=\textwidth,height=0.30\textheight,keepaspectratio]{chapter1_two_lines_accumulator.png}
		\caption{Аккумуляторный массив}
	\end{subfigure}
	\caption{Пример преобразования Хафа для изображения с двумя прямыми $y=2x-1$ и $y=-x+101$}
	\label{fig:accumulator}
\end{figure}

Если после шумоподавления в изображении сохраняется много ложных активных
точек, то в пространстве параметров возникают дополнительные максимумы, и
точность итоговых оценок прямой ухудшается. Поэтому качество предобработки
непосредственно влияет на качество оценивания параметров.

\section{Метод SSA и его комплекснозначный вариант}

Метод анализа сингулярного спектра SSA предназначен для представления ряда в
виде суммы интерпретируемых компонент, например тренда, периодических
составляющих и шума \cite{ts_structure,ssa_with_r}. Комплекснозначный вариант
Complex SSA обобщает эту конструкцию на случай комплексных рядов и особенно
удобен в задачах, где после преобразования Фурье сигнал естественно становится
комплекснозначным \cite{zhornikova2016}.

Рассмотрим комплексный временной ряд
\[
z = (z_1,\dots,z_N), \qquad z_k \in \mathbb{C}.
\]
Алгоритм CSSA состоит из четырех основных этапов.

\textbf{1. Вложение.}
Выбирается длина окна $L$, где $1 < L < N$, и строится траекторная матрица
\[
\mathbf{X} = \mathcal{T}_L(z) =
\begin{pmatrix}
	z_1 & z_2 & \dots & z_K \\
	z_2 & z_3 & \dots & z_{K+1} \\
	\vdots & \vdots & \ddots & \vdots \\
	z_L & z_{L+1} & \dots & z_N
\end{pmatrix},
\qquad K = N-L+1.
\]
Эта матрица ганкелева, то есть элементы на ее побочных диагоналях совпадают.

\textbf{2. Сингулярное разложение.}
Для траекторной матрицы строится матрица
\[
\mathbf{C}_{\mathbf{X}} = \mathbf{X}\mathbf{X}^{*},
\]
где $\mathbf{X}^{*} = \overline{\mathbf{X}}^{T}$ --- эрмитово сопряженная
матрица, где $\overline{A}$ обозначает комплексное сопряжение элементов
произвольной матрицы $A$. Матрица $\mathbf{C}_{\mathbf{X}}$ является эрмитовой
и неотрицательно определенной. Поэтому все ее собственные числа неотрицательны
и могут быть упорядочены по невозрастанию:
\[
\lambda_1 \ge \lambda_2 \ge \dots \ge \lambda_L \ge 0.
\]
Если $d=\max\{i\colon \lambda_i>0\}$, то сингулярное разложение записывается
в виде
\[
\mathbf{X} = \sum_{i=1}^{d}\mathbf{X}_i
=
\sum_{i=1}^{d}\sqrt{\lambda_i}\,U_iV_i^{*},
\]
где $U_i$ и $V_i$ --- левые и правые сингулярные векторы. Числа
$\sqrt{\lambda_i}$ называются сингулярными числами. Наборы
$(\sqrt{\lambda_i},U_i,V_i)$ принято называть собственными тройками, а $d$ ---
рангом временного ряда.

\textbf{3. Группировка.}
Выделяется множество индексов $I$, соответствующее компонентам сигнала, и
строится сумма
\[
\mathbf{X}_I = \sum_{i\in I}\mathbf{X}_i.
\]
Именно на этом шаге отделяется предполагаемая часть ряда, соответствующая сигналу, от шумовой.

\textbf{4. Восстановление.}
К матрице $\mathbf{X}_I$ применяется оператор диагонального усреднения
$\mathcal{H}_L$, в результате чего получается восстановленный ряд
\[
\widehat{z} = \mathcal{T}_L^{-1}\bigl(\mathcal{H}_L(\mathbf{X}_I)\bigr).
\]

\section{Четыре варианта CSSA для изображений}

Матрица изображения
$\mathbf{M}\in\mathbb{R}^{N_{\mathrm{row}}\times N_{\mathrm{col}}}$
рассматривается как система одномерных рядов. После применения прямого
преобразования Фурье по строкам каждая строка изображения переходит в
комплексный ряд, который можно обрабатывать методом CSSA либо по строкам, либо
по столбцам. Полученная комплексная матрица затем восстанавливается с помощью
обратного преобразования Фурье по строкам или по столбцам. Выбор измерения для
CSSA и измерения для обратного преобразования порождает четыре естественных
варианта алгоритма.

Обозначим через $\mathcal{F}_{\mathrm{row}}$ прямое преобразование Фурье по
строкам матрицы, через $\mathcal{F}^{-1}_{\mathrm{row}}$ и
$\mathcal{F}^{-1}_{\mathrm{col}}$ --- обратные преобразования по строкам и
столбцам, а через $\mathcal{C}_{\mathrm{row}}^{(r)}$ и
$\mathcal{C}_{\mathrm{col}}^{(r)}$ --- применение CSSA к каждой строке или
каждому столбцу с сохранением $r$ компонент. Тогда рассматриваются следующие
четыре варианта:
\begin{align*}
\widehat{\mathbf{S}}_{\mathrm{ROW.ROW}}
&=
\operatorname{Re}\!\Bigl(
\mathcal{F}^{-1}_{\mathrm{row}}
\bigl(
\mathcal{C}_{\mathrm{row}}^{(r)}(\mathcal{F}_{\mathrm{row}}(\mathbf{M}))
\bigr)
\Bigr), \\
\widehat{\mathbf{S}}_{\mathrm{ROW.COL}}
&=
\operatorname{Re}\!\Bigl(
\mathcal{F}^{-1}_{\mathrm{col}}
\bigl(
\mathcal{C}_{\mathrm{row}}^{(r)}(\mathcal{F}_{\mathrm{row}}(\mathbf{M}))
\bigr)
\Bigr), \\
\widehat{\mathbf{S}}_{\mathrm{COL.ROW}}
&=
\operatorname{Re}\!\Bigl(
\mathcal{F}^{-1}_{\mathrm{row}}
\bigl(
\mathcal{C}_{\mathrm{col}}^{(r)}(\mathcal{F}_{\mathrm{row}}(\mathbf{M}))
\bigr)
\Bigr), \\
\widehat{\mathbf{S}}_{\mathrm{COL.COL}}
&=
\operatorname{Re}\!\Bigl(
\mathcal{F}^{-1}_{\mathrm{col}}
\bigl(
\mathcal{C}_{\mathrm{col}}^{(r)}(\mathcal{F}_{\mathrm{row}}(\mathbf{M}))
\bigr)
\Bigr).
\end{align*}
Из этих четырех вариантов далее рассматриваются \texttt{ROW.ROW} и
\texttt{COL.ROW}. Для вариантов \texttt{ROW.COL} и \texttt{COL.COL} уже в
случае одной прямой возникают сильные ограничения на вид изображения
\cite{zhornikova2016}.

\begin{statement}[\cite{zhornikova2016}]
Пусть матрица сигнала $\mathbf{S}\in\mathbb{R}^{N_{\mathrm{row}}\times N_{\mathrm{col}}}$
содержит одну прямую
\[
S_{ij}=c\,\mathbf{1}\{i=aj+b\}, \qquad c\ne 0,
\]
и пусть в алгоритмах CSSA сохраняется одна компонента. Тогда:
\begin{enumerate}
	\item вариант \texttt{ROW.COL} точно восстанавливает $\mathbf{S}$ только в частном случае
	$a=1$, $b=0$ и $N_{\mathrm{row}}=N_{\mathrm{col}}$;
	\item вариант \texttt{COL.COL} точно восстанавливает $\mathbf{S}$ только в частном случае
	$a=1$, $b=0$ и $N_{\mathrm{row}}=N_{\mathrm{col}}$.
\end{enumerate}
\end{statement}

Так как условия $a=1$, $b=0$ и $N_{\mathrm{row}}=N_{\mathrm{col}}$ являются
слишком ограничительными для рассматриваемых изображений, варианты
\texttt{ROW.COL} и \texttt{COL.COL} далее не используются.

Теоретическое обоснование вариантов \texttt{ROW.ROW} и \texttt{COL.ROW},
включая точное восстановление в случае одной прямой и оценку ранга для
нескольких прямых, приведено во второй главе.

Ниже приведено алгоритмическое описание вариантов \texttt{ROW.ROW} и
\texttt{COL.ROW}, которые будут использоваться в дальнейших разделах.

\begin{algorithm}[H]
	\caption{Алгоритм \texttt{CSSA ROW.ROW}}
	\begin{algorithmic}[1]
		\Require матрица изображения $\mathbf{M}\in\mathbb{R}^{N_{\mathrm{row}}\times N_{\mathrm{col}}}$, число сохраняемых компонент $r$
		\Ensure восстановленное изображение $\widehat{\mathbf{S}}$
		\For{$p=1,\dots,N_{\mathrm{row}}$}
			\State вычислить DFT строки: $\boldsymbol{\xi}^{(p)} \gets \mathcal{F}(\mathbf{m}^{(p)})$
			\State применить CSSA к ряду $\boldsymbol{\xi}^{(p)}$
			\State сохранить $r$ компонент и получить восстановленный ряд $\widehat{\boldsymbol{\xi}}^{(p)}$
			\State вычислить обратное DFT:
			$\widehat{\mathbf{s}}^{(p)} \gets \operatorname{Re}(\mathcal{F}^{-1}(\widehat{\boldsymbol{\xi}}^{(p)}))$
		\EndFor
		\State составить матрицу $\widehat{\mathbf{S}}$ из строк $\widehat{\mathbf{s}}^{(1)},\dots,\widehat{\mathbf{s}}^{(N_{\mathrm{row}})}$
		\State \Return $\widehat{\mathbf{S}}$
	\end{algorithmic}
\end{algorithm}

\begin{algorithm}[H]
	\caption{Алгоритм \texttt{CSSA COL.ROW}}
	\begin{algorithmic}[1]
		\Require матрица изображения $\mathbf{M}\in\mathbb{R}^{N_{\mathrm{row}}\times N_{\mathrm{col}}}$, число сохраняемых компонент $r$
		\Ensure восстановленное изображение $\widehat{\mathbf{S}}$
		\State $\mathbf{Y}_{p,\cdot}\gets \mathcal{F}(\mathbf{m}^{(p)})$, $p=1,\dots,N_{\mathrm{row}}$
		\For{$s=1,\dots,N_{\mathrm{col}}$}
			\State $\widehat{\mathbf{y}}_{s}\gets$ результат CSSA для $\mathbf{Y}_{\cdot,s}$ с сохранением $r$ компонент
		\EndFor
		\State $\widehat{\mathbf{Y}} \gets (\widehat{\mathbf{y}}_{1},\dots,\widehat{\mathbf{y}}_{N_{\mathrm{col}}})$
		\State $\widehat{\mathbf{s}}^{(p)} \gets \operatorname{Re}(\mathcal{F}^{-1}(\widehat{\mathbf{Y}}_{p,\cdot}))$, $p=1,\dots,N_{\mathrm{row}}$
		\State $\widehat{\mathbf{S}} \gets (\widehat{\mathbf{s}}^{(1)},\dots,\widehat{\mathbf{s}}^{(N_{\mathrm{row}})})^{\top}$
		\State \Return $\widehat{\mathbf{S}}$
	\end{algorithmic}
\end{algorithm}

На одном и том же зашумленном изображении варианты \texttt{ROW.ROW} и
\texttt{COL.ROW} дают разные восстановленные матрицы: первый метод
обрабатывает строки независимо, а второй использует структуру столбцов матрицы
после DFT по строкам. Типичный результат работы этих двух вариантов показан на
рис.~\ref{fig:cssa-row-col-examples}. Здесь синим цветом выделены элементы
матрицы с отрицательными значениями.

\begin{figure}[H]
	\centering
	\begin{subfigure}{0.48\textwidth}
		\centering
		\includegraphics[width=\textwidth]{chapter1_cssa_row_row_processed.png}
		\caption{\texttt{CSSA ROW.ROW}}
	\end{subfigure}
	\hfill
	\begin{subfigure}{0.48\textwidth}
		\centering
		\includegraphics[width=\textwidth]{chapter1_cssa_col_row_processed.png}
		\caption{\texttt{CSSA COL.ROW}}
	\end{subfigure}
	\caption{Пример предобработки изображения двумя вариантами CSSA без последующей пороговой фильтрации}
	\label{fig:cssa-row-col-examples}
\end{figure}

\section{Стандартные методы предобработки}

Помимо алгоритмов на основе CSSA, в работе рассматриваются стандартные методы
шумоподавления, применяемые к той же зашумленной матрице изображения перед
использованием метода Хафа.

\subsection*{Медианный фильтр}

Медианный фильтр заменяет значение каждого пикселя медианой значений в окне
\cite{huang1979,mafi2019}. Если через $W_{ij}$ обозначить квадратное окно вокруг
пикселя $(i,j)$ размера $k$, то результат фильтрации имеет вид
\[
\widehat{M}_{ij} = \operatorname{median}\{M_{uv} \colon (u,v)\in W_{ij}\}.
\]
По умолчанию используется окно размера \(3\times 3\).
Существенное свойство медианного фильтра состоит в способности подавлять
одиночные выбросы, не размывая резкие границы столь же сильно, как линейные
сглаживающие фильтры.

\subsection*{Фильтр Винера}

Фильтр Винера в общем виде строит линейную оценку сигнала, минимизирующую
среднеквадратичную ошибку в каждом элементе матрицы \cite{wiener1949}. Далее
рассматривается локально-адаптивная модификация этого подхода для изображений;
соответствующее описание приводится, например, у Lim \cite{lim1990}, а сравнение
с медианным фильтром в задаче шумоподавления изображений приведено в
\cite{mafi2019}. Термин
«локальный» означает, что параметры фильтра, как и в случае медианного фильтра,
определяются отдельно для каждого пикселя по его окну.

Пусть $W_{ij}$ --- квадратное окно размера $k\times k$ с центром в точке
$(i,j)$. Стандартным считается окно размера \(5\times 5\). Тогда для каждого
пикселя сначала вычисляются локальное среднее
\[
\mu_{ij} = \frac{1}{|W_{ij}|}\sum_{(u,v)\in W_{ij}} M_{uv}
\]
и локальная дисперсия
\[
\nu_{ij} =
\frac{1}{|W_{ij}|}\sum_{(u,v)\in W_{ij}} M_{uv}^2 - \mu_{ij}^2.
\]
После этого вводится величина $\sigma_n^2$, обозначающая оценку дисперсии
шумовой компоненты изображения. В рассматриваемой схеме эта величина
определяется формулой
\[
\sigma_n^2 = \frac{1}{N_{\mathrm{row}}N_{\mathrm{col}}}\sum_{i=1}^{N_{\mathrm{row}}}\sum_{j=1}^{N_{\mathrm{col}}} \nu_{ij},
\]
то есть как среднее значение локальных дисперсий по всем пикселям изображения. Далее результат записывается как
\[
\widehat{M}_{ij}
=
\mu_{ij}
+
\frac{\max(\nu_{ij}-\sigma_n^2,\,0)}{\max(\nu_{ij},\,\varepsilon)}
\left(M_{ij}-\mu_{ij}\right).
\]
Здесь $\varepsilon>0$ --- малая константа, предотвращающая деление на ноль.
Приведенная формула показывает, что при малой локальной дисперсии, сравнимой с
оценкой шума, поправочный множитель оказывается малым, и значение пикселя
смещается к локальному среднему. Если же локальная дисперсия велика, степень
сглаживания уменьшается, что способствует сохранению локальной структуры
изображения. Таким образом, фильтр согласуется с локальной структурой
изображения: в почти однородных областях он подавляет шум сильнее, а в областях
с заметной вариативностью в большей степени сохраняет полезный сигнал.

\subsection*{\texttt{MAX.ROW}}

Для изображений, состоящих из небольшого числа прямых, естественна простая
строковая эвристика: в каждой строке оставить только несколько наибольших
значений, а остальные занулить. Если предполагается, что в строке присутствует
не более $k$ пикселей сигнала, можно определить оператор
\[
\mathcal{M}_{k}(\mathbf{M})_{ij}
=
\begin{cases}
M_{ij}, & j \in J_i^{(k)},\\
0, & j \notin J_i^{(k)},
\end{cases}
\]
где $J_i^{(k)}$ --- множество индексов $k$ наибольших элементов в $i$-й
строке. После этого преобразование Хафа применяется к матрице
$\mathcal{M}_{k}(\mathbf{M})$.

В дальнейшем этот подход следует сравнивать с \texttt{CSSA ROW.ROW} и
\texttt{CSSA COL.ROW}. Он может быть полезен, когда пиксели сигнала
действительно являются крупнейшими значениями в строке, но может работать
некорректно при сильном шуме, пересечениях прямых и разных интенсивностях
линий.

На рис.~\ref{fig:preprocessing-five-methods} приведен пример матриц,
полученных после разных вариантов предобработки одного и того же зашумленного
изображения.

\begin{figure}[H]
	\centering
	\begin{subfigure}{0.32\textwidth}
		\centering
		\includegraphics[width=\textwidth]{chapter1_cssa_row_row_processed.png}
		\caption{\texttt{CSSA ROW.ROW}}
	\end{subfigure}
	\hfill
	\begin{subfigure}{0.32\textwidth}
		\centering
		\includegraphics[width=\textwidth]{chapter1_cssa_col_row_processed.png}
		\caption{\texttt{CSSA COL.ROW}}
	\end{subfigure}
	\hfill
	\begin{subfigure}{0.32\textwidth}
		\centering
		\includegraphics[width=\textwidth]{chapter1_max_row_processed.png}
		\caption{\texttt{MAX.ROW}}
	\end{subfigure}

	\vspace{0.4em}

	\makebox[\textwidth][c]{%
		\begin{subfigure}{0.32\textwidth}
			\centering
			\includegraphics[width=\textwidth]{chapter1_median_processed.png}
			\caption{Median}
		\end{subfigure}
		\hspace{0.04\textwidth}
		\begin{subfigure}{0.32\textwidth}
			\centering
			\includegraphics[width=\textwidth]{chapter1_wiener_processed.png}
			\caption{Wiener}
		\end{subfigure}%
	}
	\caption{Результаты предобработки без пороговой фильтрации для пяти сравниваемых методов}
	\label{fig:preprocessing-five-methods}
\end{figure}

\section{Схема оценивания}

Во всех экспериментах используется одна и та же схема оценивания:
\begin{enumerate}
	\item выбирается способ формирования изображения для преобразования Хафа: стандартный фильтр, один из вариантов CSSA;
	\item по полученному изображению формируется множество активных элементов матрицы;
	\item по этому множеству строится аккумуляторный массив метода Хафа;
	\item по максимумам аккумулятора находятся оценки $(\hat{\rho},\hat{\theta})$.
\end{enumerate}

В дальнейшем качество методов сравнивается по среднеквадратическим ошибкам
оценок параметров прямой. Если проводится $R$ независимых повторений
эксперимента и истинные параметры равны $(\rho,\theta)$, то для полученных
оценок $(\widehat{\rho}^{(r)},\widehat{\theta}^{(r)})$, $r=1,\dots,R$,
используются метрики
\[
\operatorname{MSE}(\widehat{\rho})
=
\frac{1}{R}\sum_{r=1}^{R}
\bigl(\widehat{\rho}^{(r)}-\rho\bigr)^2,
\qquad
\operatorname{MSE}(\widehat{\theta})
=
\frac{1}{R}\sum_{r=1}^{R}
\bigl(\widehat{\theta}^{(r)}-\theta\bigr)^2.
\]

